% page 557 of the facsimile (image p-556 of the scan)
% block 177.8 x 263.6 mm ; left margin 24.2 mm
\pgc{2.540mm}{0.000mm}{
\l{}
\l{}
\l{}
\l{}
\l{\cel{57}549}
\l{}
\l{\cel{3}\sou{squilo}. (\sou{bot.)}\cel{2}Genero de \textquotedbl{}liliacei\textquotedbl{}. L.Scilla maritima.}
\l{}
\l{\cel{3}st !\cel{3}Vorto quan onu dicas ad ulu por ke lu tacez.}
\l{}
\l{\cel{3}\sou{stabo}.\cel{2}(\sou{armeo)}. Ensemblo de oficiri quin onu selektis por servar}
\l{\cel{6}oficiro superiora qua komandas armeo, ed esar apud lu quaze}
\l{\cel{6}kom konsilantaro qua transmisas ed obediigas lua komandi. -}
\l{\cel{6}DER.}
\l{}
\l{\cel{3}\sou{stabila}. Qua tendencas permanar en la sama loko, en la\cel{1}\sou{sama}\cel{1}si-}
\l{\cel{6}tueso, en la sama stando. -\cel{2}DEFIS.}
\l{}
\l{\cel{3}\sou{stablo}.\cel{2}Klozajo muroza e tektoza ube onu lojigas la bestii, la}
\l{\cel{6}bruti. - DEFIS.}
\l{}
\l{\cel{3}\sou{stacar}. (\sou{netrans.}) (\sou{aludante homo}). Havar sua korpo vertikala\cel{1}\sou{sur}}
\l{\cel{6}sua pedi, t.e. standar ne-sidante, ne-kushinte, ne-jacante,ne}
\l{\cel{6}genu-pozinte, e ne-inklinite. - I.}
\l{}
\l{\cel{3}\sou{staciono}. I. Singla ek la loki ube haltas treno, por embarkar o}
\l{\cel{6}desembarkar pasajeri o vari. - II. Loko ube esas, o haltas,}
\l{\cel{6}la veturi publik-uzebla (autofiakri - F.\cel{1}\sou{taxi}\cel{1}- autobusi,}
\l{\cel{6}e c.).- DEFIS.}
\l{}
\l{\cel{3}\sou{stacionara}. Qua haltigas sua movo e permanas dum kelka tempo en}
\l{\cel{6}la sama loko (o stando). - DEFIS.}
\l{}
\l{\cel{3}\sou{stadio}. (\sou{antique}). Areo klozita, longa ye 180 metri, ube onu}
\l{\cel{6}interluktis por obtenar la premio pri homokuro. - (\sou{nun)}.}
\l{\cel{6}\sou{Lico}\cel{1}por la konkursi atletala, automobilala, biciklala. -}
\l{\cel{6}DEFIRS.}
\l{}
\l{\cel{3}\sou{stafilo}. (\sou{bot.)}\cel{2}Arboreto, vicina ye la aceri, kun folii kompozita,}
\l{\cel{6}kun aristi plumatra. -L.\cel{1}\sou{staphylea}.}
\l{}
\l{\cel{3}\sou{stafisagro}. (\sou{bot.})\cel{2}Planto ek la familio \textquotedbl{}ranunkulacei\textquotedbl{}. L.\cel{1}\sou{Del}-}
\l{\cel{6}\sou{phinium staphis agria}.}
\l{}
\l{\cel{3}\sou{stagnar}. (\sou{netrans.})\cel{2}Qua ne fluas, qua ne cirkulas, que ne rinov-}
\l{\cel{6}eskas e, pro lo, eventuale putreskas. - (\sou{anke metaf.}) - DEFIS.}
\l{}
\l{\cel{3}\sou{stakiso.}\cel{2}(\sou{bot.})\cel{2}Planto ek la familio \textquotedbl{}labiei\textquotedbl{}, deveninta de}
\l{\cel{6}Japonia,di qua la+rizomi esas manjebla ed analoga a salsifio.}
\l{\cel{6}- S.\cel{1}\sou{stachys}. -}
\l{}
\l{\cel{3}\sou{stalo}.\cel{2}Metalo ek fero pura a qua esas enkorpigita ula quanto de}
\l{\cel{6}karbo - plu harda, e di qua la grani esas plu tenua e plu}
\l{\cel{6}briloza kam olti di fero. - DER.}
\l{}
\l{\cel{3}\sou{stalagmito}. Agreguro, sur la sulo di kavajo subtera, per infil-}
\l{\cel{6}truri qui falas de la vulto. - DEFIRS.}
\l{}
\l{\cel{3}stalaktito. Agreguro an la vulto di kavajo subtera, per infil-}
\l{\cel{6}truri qui kontenas dissolvite sali kalkoza, silicoza, e c.}
\l{\cel{61}DEFIRS.}
}
